\documentclass[11pt,a4paper]{article}

\usepackage{amsmath,amssymb,amsthm}
\usepackage{hyperref}
\usepackage[margin=1in]{geometry}
\usepackage{xcolor}
\usepackage{booktabs}
\usepackage{array}
\usepackage{float}
\usepackage{mathtools}

\newtheorem{theorem}{Theorem}[section]
\newtheorem{conjecture}[theorem]{Conjecture}
\newtheorem{corollary}[theorem]{Corollary}
\newtheorem{lemma}[theorem]{Lemma}
\newtheorem{proposition}[theorem]{Proposition}
\newtheorem{definition}[theorem]{Definition}
\theoremstyle{remark}
\newtheorem{remark}[theorem]{Remark}

\newcommand{\PT}{\mathbb{PT}}
\newcommand{\CP}{\mathbb{CP}}
\newcommand{\C}{\mathbb{C}}
\newcommand{\Z}{\mathbb{Z}}
\newcommand{\R}{\mathbb{R}}
\newcommand{\F}{\mathbb{F}}
\newcommand{\cO}{\mathcal{O}}
\newcommand{\cC}{\mathcal{C}}
\newcommand{\cN}{\mathcal{N}}
\newcommand{\cU}{\mathcal{U}}
\newcommand{\cW}{\mathcal{W}}
\newcommand{\cG}{\mathcal{G}}
\newcommand{\cF}{\mathcal{F}}
\newcommand{\PP}{\mathcal{P}}
\newcommand{\barPP}{\overline{\mathcal{P}}}


\title{The Equiangular Characterization of Mermin Pentagrams:\\
       Symplectic Embedding, 10-Ray Cap, and the Failure\\
       of the Collinearity--Obstruction Dictionary\\
       {\large (Paper X)}}


\author{Zhou-Li Chen \\ co-nlang Research}
\date{May 2026}

\begin{document}

\maketitle

\begin{abstract}
  Paper~IX~\cite{PaperIX} established the Fano--pentagram correspondence
  (Proposition~2.6) and conjectured a ``collinearity--obstruction
  dictionary'' (Conjecture~2.5): the Kochen--Specker obstruction
  $[f_3] \in H^1(K_5, \cF)$ would be equivalent to a collinearity
  obstruction in $PG(2, \F_2)$.

  Paper~X tests this conjecture computationally and finds it false.
  Among 12{,}096 Mermin pentagrams in $W(5, \F_2)$, we classify them
  into four types by their Fano-point shadow:
  Type~A (5 distinct points) $= 0$,
  Type~B (4 distinct) $= 1{,}344$,
  Type~F2 (2 transverse) $= 7{,}168$,
  Type~F4 (4 transverse) $= 3{,}584$.
  Collinearity is a weak statistical signal, not a deterministic
  constraint: 33.3\% of Type~B are 4-arcs with no collinearity, and
  25\% of Type~F4 have \emph{zero} Fano points at all.

  The true invariant is geometric, not projective. We prove
  (Theorem~\ref{thm:equiangular}) that in the proper context framework
  of $W(5, \F_2)$, three notions coincide:
  \[
  K_5 \;\Longleftrightarrow\; \text{equiangular } K_5
  \;\Longleftrightarrow\; \text{Mermin pentagram}.
  \]
  The equiangular condition $\dim(L_i \cap L_j) = 1$ for all 10 pairs
  is \emph{forced} by the $K_5$ structure itself (pure geometry, no
  computation needed). The 10 intersection rays form a 10-point cap in
  $PG(5, \F_2)$ (Theorem~\ref{thm:cap}).

  The odd parity of every equiangular $K_5$ is the evaluation of
  $[f_3] \in H^1(K_5, \cF)$, a symplectic characteristic class of
  $W(5, \F_2)$. Conjecture~2.5 asked the wrong question: collinearity
  is a projection of the equiangular structure onto $V$, not the
  obstruction itself.
\end{abstract}

%------------------------------------------------------------------
\section{Introduction}
\label{sec:intro}
%------------------------------------------------------------------

Paper~IX~\cite{PaperIX} embedded the Fano plane $PG(2, \F_2)$ as an
isotropic subgeometry of $W(5, \F_2)$ via the hyperbolic embedding
$GL(3, \F_2) \hookrightarrow PSp(6, \F_2)$, and established a
surjective map from Mermin pentagram contexts to Fano points
(Proposition~2.6). Conjecture~2.5 proposed a ``collinearity--obstruction
dictionary'': the KS obstruction $[f_3] \in H^1(K_5, \cF)$ would
translate to a collinearity obstruction in $PG(2, \F_2)$.

This paper tests Conjecture~2.5 by exhaustive computation and
geometric analysis. The main results are:

\begin{enumerate}
  \item \textbf{Theorem~\ref{thm:type-a} (Type~A Exclusion, computational).}
    No Mermin pentagram has 5 distinct Fano points forming 2 intersecting
    Fano lines. The sign algebra alone does not rule out Type~A; the
    obstruction is symplectic-geometric.

  \item \textbf{Theorem~\ref{thm:equiangular} (Equiangular Characterization, geometric).}
    In the proper context framework of $W(5, \F_2)$:
    \[
    K_5 \;\Longleftrightarrow\; \text{equiangular } K_5
    \;\Longleftrightarrow\; \text{Mermin pentagram}.
    \]
    Non-equiangular $K_5$ do not exist. Non-Mermin equiangular $K_5$
    do not exist.

  \item \textbf{Theorem~\ref{thm:cap} (10-Ray Cap Structure).}
    The 10 intersection rays $R_{ij} = L_i \cap L_j$ form a 10-point
    cap in $PG(5, \F_2)$: no three rays are collinear. Each Lagrangian
    contains 4 rays summing to zero. The KS contradiction arises from
    5 Lagrangian dependencies summing to $0 = 1 \pmod{2}$.

  \item \textbf{Corollary~\ref{cor:characteristic} ($[f_3]$ as Symplectic
    Characteristic Class).} The odd parity of every equiangular $K_5$
    is the evaluation of $[f_3] \in H^1(K_5, \cF)$, a global
    cohomological invariant of $W(5, \F_2)$.
\end{enumerate}

The core narrative:

\begin{quote}
  Conjecture~2.5 asked the wrong question. Collinearity is a
  \emph{shadow} --- a projection of the equiangular Lagrangian
  configuration onto the Fano subgeometry $V$. The true invariant is
  the symplectic embedding of $K_5$ into $W(5, \F_2)$: every proper
  context $K_5$ is equiangular ($\dim(L_i \cap L_j) = 1$ for all
  10 pairs), and this rigid geometric structure forces the KS parity
  anomaly via the symplectic characteristic class $[f_3]$.
\end{quote}

%------------------------------------------------------------------
\section{Computational Enumeration}
\label{sec:enumeration}
%------------------------------------------------------------------

\subsection{Lagrangians and Proper Contexts}

The symplectic polar space $W(5, \F_2)$ has 63 points and 315 lines.
A \emph{Lagrangian} is a 3-dimensional totally isotropic subspace
(maximal isotropic). There are exactly 135 Lagrangians in $W(5, \F_2)$.

Each Lagrangian $L$ has 7 non-zero points, forming a Fano plane
$PG(2, \F_2)$. A \emph{proper context} is a 4-element subset of $L$
obtained by removing a Fano line (3 collinear points). Each Lagrangian
has 7 Fano lines, hence 7 proper contexts.

\begin{table}[h]
\centering
\begin{tabular}{l r l}
\toprule
Object & Count & Notes \\
\midrule
Lagrangians in $W(5, \F_2)$ & 135 & 3-dim totally isotropic subspaces \\
Proper contexts & 945 & $135 \times 7$ (complement of each Fano line) \\
Contexts with $-\mathbf{I}$ product & 324 & 34.3\% of 945 \\
\bottomrule
\end{tabular}
\caption{Enumeration of Lagrangians and proper contexts.}
\label{tab:enumeration}
\end{table}

\subsection{\texorpdfstring{MASA Classification by $V$-Intersection}{MASA Classification by V-Intersection}}

Under the hyperbolic embedding $GL(3, \F_2) \hookrightarrow PSp(6, \F_2)$,
the Lagrangian $V \subset V \oplus V^*$ is the Fano isotropic
subgeometry. Each Lagrangian $L$ intersects $V$ in a subspace of
dimension $k = \dim(L \cap V)$:

\begin{table}[h]
\centering
\begin{tabular}{c c c l}
\toprule
Type & $\dim(L \cap V)$ & $|L \cap V \setminus \{0\}|$ & Count \\
\midrule
$k=7$ & 3 & 7 & 1 ($V$ itself) \\
$k=3$ & 2 & 3 & 14 (contains a Fano line) \\
$k=1$ & 1 & 1 & 56 (contains a single Fano point) \\
$k=0$ & 0 & 0 & 64 (transverse to $V$) \\
\bottomrule
\end{tabular}
\caption{MASA classification by intersection with $V$.}
\label{tab:masa-types}
\end{table}

\subsection{Mermin Pentagrams}

A \emph{Mermin pentagram}~\cite{Mermin1990} is a $K_5$ clique of proper contexts:
5 contexts $\{C_1, \ldots, C_5\}$ with 10 distinct shared operators
(each operator in exactly 2 contexts), satisfying the odd-parity
condition $\prod_{i=1}^5 s_i = -1$ where $s_i = \pm 1$ is the product
sign of context $C_i$.

Exhaustive enumeration yields:

\begin{table}[h]
\centering
\begin{tabular}{l r l}
\toprule
Object & Count & Notes \\
\midrule
Mermin pentagrams & \textbf{12{,}096} & verified by $G_2(2)$ orbit decomposition \\
$G_2(2)$ orbits & 4 & sizes $6{,}048 + 2{,}016 + 2{,}016 + 2{,}016$ \\
\bottomrule
\end{tabular}
\caption{Mermin pentagram enumeration.}
\label{tab:mermin-count}
\end{table}

\begin{remark}[Previous Overcounting]
An earlier version of the enumeration script counted each $K_5$ once
per root vertex, yielding $5 \times 12{,}096 = 60{,}480$. After
deduplication, the correct count is $12{,}096$, matching the
literature value for the $G_2(2)$ orbit of the standard pentagram.
\end{remark}

\subsection{Type Classification by Fano Shadow}

Each Mermin pentagram has 5 contexts, each contributing a Fano point
(if the context's Lagrangian intersects $V$ non-trivially) or being
\emph{transverse} (if the Lagrangian is transverse to $V$). We classify
by the number of distinct Fano points:

\begin{table}[h]
\centering
\begin{tabular}{l l r r}
\toprule
Type & Pattern & Count & Fraction \\
\midrule
\textbf{A} & 5 distinct Fano points & \textbf{0} & \textbf{0\%} \\
B & 4 distinct (one collision) & 1{,}344 & 11.11\% \\
F2 & 3 Fano points + 2 transverse & 7{,}168 & 59.26\% \\
F4 & 1 Fano point + 4 transverse & 3{,}584 & 29.63\% \\
\midrule
Total & & \textbf{12{,}096} & 100\% \\
\bottomrule
\end{tabular}
\caption{Mermin pentagram type classification.}
\label{tab:type-classification}
\end{table}

%------------------------------------------------------------------
\section{Type~A Exclusion}
\label{sec:type-a}
%------------------------------------------------------------------

\begin{theorem}[Type~A Exclusion]
\label{thm:type-a}
No Type~A Mermin pentagram exists. Among all $K_5$ context-quintets
with 5 distinct Fano points forming 2 intersecting Fano lines, none
satisfy the KS odd-parity condition.
\end{theorem}

\begin{proof}[Proof (computational)]
Exhaustive enumeration over all 945 proper contexts confirms that
zero $K_5$ cliques with 5 distinct Fano points satisfy the odd-parity
condition. The code is available in the supplementary
material~\cite{PaperXsuppl}.
\end{proof}

\begin{remark}[Why the Naive Sign Argument Fails]
Let the 5 Fano points form two intersecting Fano lines
$L_1 = \{p_1, p_2, p_3\}$ and $L_2 = \{p_3, p_4, p_5\}$ sharing $p_3$.
Collinearity constraints give:
\[
s_1 \cdot s_2 \cdot s_3 = -1 \quad (\text{from } L_1), \qquad
s_3 \cdot s_4 \cdot s_5 = -1 \quad (\text{from } L_2).
\]
These are compatible with $s_3 = -1$: set $s_1 s_2 = +1$,
$s_4 s_5 = +1$, then total product $= s_3 = -1$ (Mermin $\checkmark$).
The sign algebra alone does \emph{not} rule out Type~A.

The real obstruction is symplectic-geometric: the $K_5$ structure in
$W(5, \F_2)$ cannot realize the $s_3 = -1$ sign combination for Type~A
configurations. A geometric proof (without computation) is open.
\end{remark}

\begin{corollary}
\label{cor:type-a}
Conjecture~2.5 cannot be tested on Type~A pentagrams. The
collinearity--obstruction dictionary must be reformulated for Types~B,
F2, or F4.
\end{corollary}

%------------------------------------------------------------------
\section{Collinearity is Not the Obstruction}
\label{sec:collinearity-failure}
%------------------------------------------------------------------

\subsection{Type~B: 4-Arcs and Collinear Triples}

Type~B has 4 distinct Fano points with one collision (2 contexts share
the same Fano point). Among 1{,}344 Type~B pentagrams:

\begin{itemize}
  \item 896 (66.7\%) have 1 collinear triple among the 4 Fano points
  \item 448 (33.3\%) are \emph{4-arcs}: no 3 points collinear
\end{itemize}

The existence of 4-arcs (hyperovals in $PG(2, \F_2)$) shows that valid
KS pentagrams can have \emph{no collinearity structure at all}.

For the 896 collinear cases, the sign pattern of the collinear
triple is:

\begin{table}[h]
\centering
\begin{tabular}{l r r}
\toprule
\# of $-I$ in collinear triple & Count & Fraction \\
\midrule
0 $(+,+,+)$ & 300 & 33.5\% \\
1 $(-,+,+)$ + perms & 444 & 49.6\% \\
2 $(-,-,+)$ + perms & 132 & 14.7\% \\
3 $(-,-,-)$ & 20 & 2.2\% \\
\bottomrule
\end{tabular}
\caption{Type~B collinear triple sign distribution.}
\label{tab:type-b-signs}
\end{table}

The naive expectation ``collinear $\Rightarrow$ product $= -1$'' is
\textbf{false}: the distribution is approximately 50--50, not
concentrated on odd parity.

\subsection{Type~F2: The Dominant Case}

Type~F2 has 3 Fano points (from 3 non-transverse contexts) plus
2 transverse contexts. Among 7{,}168 Type~F2 pentagrams:

\begin{itemize}
  \item 896 (12.5\%) have 3 collinear Fano points
  \item 6{,}272 (87.5\%) have non-collinear Fano points
\end{itemize}

The non-transverse context sign distribution (all Type~F2):

\begin{table}[h]
\centering
\begin{tabular}{l r r}
\toprule
\# of $-I$ in 3 non-transverse & Count & Fraction \\
\midrule
0 $(+,+,+)$ & 3{,}028 & 42.2\% \\
1 $(-,+,+)$ + perms & 3{,}036 & 42.4\% \\
2 $(-,-,+)$ + perms & 972 & 13.6\% \\
3 $(-,-,-)$ & 132 & 1.8\% \\
\bottomrule
\end{tabular}
\caption{Type~F2 sign distribution.}
\label{tab:type-f2-signs}
\end{table}

The collinear subset has a slightly different distribution (more
1-minus: 46.9\% vs 42.4\%), but the effect is modest. Collinearity
is a weak statistical signal, not a deterministic constraint.

\subsection{Summary: Collinearity is a Shadow}

\begin{proposition}[Naive Dictionary Failure]
\label{prop:naive-failure}
The collinearity pattern of Fano points does \emph{not} determine the
KS obstruction:
\begin{enumerate}
  \item 33.3\% of Type~B are 4-arcs with no collinearity.
  \item 87.5\% of Type~F2 have non-collinear Fano points.
  \item Collinear sign patterns are approximately 50--50, not
    concentrated on odd parity.
\end{enumerate}
Collinearity is the KS obstruction's \emph{shadow} on $V$, not the
obstruction itself.
\end{proposition}

%------------------------------------------------------------------
\section{Type~F4: Pure Transverse Obstruction}
\label{sec:type-f4}
%------------------------------------------------------------------

\subsection{Zero Fano Points}

Type~F4 has 4 transverse contexts plus 1 context that may or may not
contribute a Fano point. Among 3{,}584 Type~F4 pentagrams:

\begin{itemize}
  \item 2{,}688 (75\%) have 1 Fano point, uniformly distributed
    (384 each for all 7 Fano points)
  \item \textbf{896 (25\%) have ZERO Fano points} --- all 5 contexts
    are transverse
\end{itemize}

The 896 ``pure transverse'' pentagrams are the strongest evidence
against the naive dictionary: they have \emph{no Fano points at all},
yet the KS obstruction is fully present.

\begin{theorem}[Pivot Pentagon Rigidity]
\label{thm:ghost}
All 896 pure transverse Type~F4 pentagrams share the same rigid
geometric structure:
\begin{itemize}
  \item \textbf{Lagrangian intersection:} All 10 pairwise intersections
    have dimension 1 (equiangular configuration).
  \item \textbf{$V$-intersection type:} $(0,0,0,0,3)$ --- 4 Lagrangians
    transverse to $V$ ($k=0$), 1 contains a Fano line ($k=3$).
\end{itemize}
\end{theorem}

\begin{remark}[Fano Geometry at the Lagrangian Level]
``Zero Fano points'' describes the \emph{context} level, not the
\emph{Lagrangian} level. The $k=3$ Lagrangian contains a Fano line,
and this Fano line is precisely the element removed to form the proper
context. The context therefore has 0 Fano points not because the Fano
geometry is absent, but because it serves as the \emph{structural
pivot} of the construction. Pivot Pentagons are the case where Fano
geometry is most tightly embedded: the Fano line is not a shadow but
the axis around which the context is built.
\end{remark}

%------------------------------------------------------------------
\section{The Equiangular Characterization}
\label{sec:equiangular}
%------------------------------------------------------------------

\subsection{The Main Theorem}

\begin{theorem}[Equiangular Characterization]
\label{thm:equiangular}
In the proper context framework of $W(5, \F_2)$, the following three
notions coincide:
\[
K_5 \;\Longleftrightarrow\; \text{equiangular } K_5
\;\Longleftrightarrow\; \text{Mermin pentagram}.
\]
\end{theorem}

\begin{proof}[Proof (pure geometry, no computation)]
The equiangular condition is forced by the $K_5$ structure itself:

\begin{enumerate}
  \item $K_5$ has 10 operators, each appearing in exactly 2 contexts.
  \item Each pair of Lagrangians $(L_i, L_j)$ must share exactly the
    operator at their edge.
  \item If $\dim(L_i \cap L_j) = 0$: they share 0 operators ---
    impossible for a $K_5$ edge.
  \item If $\dim(L_i \cap L_j) = 2$: they share 3 Fano points (a Fano
    line). But proper context construction removes the Fano line from
    each context, so the contexts share 0 operators --- impossible for
    a $K_5$ edge.
  \item Therefore $\dim(L_i \cap L_j) = 1$ for all 10 pairs. \qedhere
\end{enumerate}
\end{proof}

\begin{corollary}
Non-equiangular $K_5$ do not exist in the proper context framework.
Non-Mermin equiangular $K_5$ do not exist either.
\end{corollary}

\begin{remark}[Computational Confirmation]
The geometric theorem is confirmed computationally:

\begin{table}[h]
\centering
\begin{tabular}{l r}
\toprule
Direction & Count \\
\midrule
Mermin $\Rightarrow$ equiangular & 12{,}096 / 12{,}096 (100\%) \\
Equiangular $\Rightarrow$ Mermin & 12{,}096 / 12{,}096 (100\%) \\
Non-equiangular $K_5$ & 0 \\
Even-parity equiangular $K_5$ & 0 \\
\bottomrule
\end{tabular}
\caption{Computational verification of Theorem~\ref{thm:equiangular}.}
\label{tab:equiangular-verify}
\end{table}

The parity distribution among equiangular $K_5$
(Supplementary Script~P8~\cite{PaperXsuppl}):

\begin{table}[h]
\centering
\begin{tabular}{l r r}
\toprule
Pattern & Count & Fraction \\
\midrule
1 minus & 7{,}884 & 65.18\% \\
3 minus & 4{,}104 & 33.93\% \\
5 minus & 108 & 0.89\% \\
\textbf{All odd} & \textbf{12{,}096} & \textbf{100\%} \\
\bottomrule
\end{tabular}
\caption{Parity distribution among equiangular $K_5$ (Script~P8~\cite{PaperXsuppl}).}
\label{tab:parity-dist}
\end{table}
\end{remark}

\subsection{The 10-Ray Cap}

\begin{theorem}[10-Ray Cap Structure]
\label{thm:cap}
The 10 intersection rays $R_{ij} = L_i \cap L_j$ are exactly the
10 shared operators. They form a \emph{10-point cap} in $PG(5, \F_2)$:
no three rays are collinear. The 10 rays span the full $\F_2^6$.
\end{theorem}

\begin{proof}[Proof (computational)]
Each Lagrangian $L_i$ contains 4 rays (one for each edge incident to
vertex $i$ in $K_5$). Each ray lies in exactly 2 Lagrangians.

The cap property and spanning property are verified computationally
(Supplementary \\Script~P7~\cite{PaperXsuppl}):
for a representative equiangular $K_5$, the 10 rays $R_{ij}$ satisfy:
\begin{itemize}
  \item \textbf{Cap:} $\binom{10}{3} = 120$ triples checked;
    0 are collinear (i.e., no $R_{mn} = R_{ij} + R_{kl}$).
  \item \textbf{Span:} $\dim\langle R_{ij} \rangle = 6$
    (the 10 rays span the full $\F_2^6$).
\end{itemize}
A purely geometric proof of the cap property from the equiangular
condition remains open (see Section~\ref{sec:open-problems}).
\end{proof}

\begin{remark}[Geometric Structure]
The $K_5$ embedding in $PG(5, \F_2)$:
\begin{itemize}
  \item 5 vertices $\to$ 5 Lagrangians (3-dim subspaces)
  \item 10 edges $\to$ 10 intersection rays (1-dim subspaces)
  \item Each Lagrangian contains 4 rays; each ray lies in 2 Lagrangians
  \item No 3 rays collinear (cap property)
\end{itemize}
\end{remark}

\subsection{The KS Contradiction from Equiangular Geometry}

The KS contradiction separates into two independent layers:
the \emph{linear algebra} of $\F_2^6$ and the \emph{sign contradiction}
on $\{\pm 1\}$.

\medskip
\noindent\textbf{Layer 1: Linear dependencies in $\F_2^6$.}

\begin{enumerate}
  \item Each proper context $C_i$ consists of 4 operators whose
    representative vectors in $\F_2^6$ sum to zero:
    $\sum_{v \in C_i} v = 0$.
    (The 4 operators form the complement of a Fano line in a
    3-dimensional Lagrangian, hence are linearly dependent.)
  \item Consequently, the Pauli product of the 4 operators in $C_i$
    is $\pm I$, with sign $s_i \in \{+1, -1\}$ determined by the
    symplectic phase.
\end{enumerate}

\medskip
\noindent\textbf{Layer 2: The sign contradiction.}

Suppose a consistent value assignment $f\colon \{\text{operators}\} \to
\{\pm 1\}$ exists.

\begin{enumerate}
  \setcounter{enumi}{2}
  \item For each context $C_i$, the product of assigned values satisfies
    $\prod_{v \in C_i} f(v) = s_i$.
  \item Taking the product over all 5 contexts:
    \[
    \prod_{i=1}^{5} \prod_{v \in C_i} f(v) = \prod_{i=1}^{5} s_i.
    \]
  \item \textbf{LHS:} Each of the 10 operators appears in exactly
    2 contexts, so $\prod_{i} \prod_{v \in C_i} f(v)
    = \prod_{v} f(v)^2 = (+1)^{10} = +1$.
  \item \textbf{RHS:} The odd-parity condition (Definition,
    \S\ref{sec:enumeration}) requires $\prod_{i=1}^{5} s_i = -1$.
  \item Hence $+1 = -1$: contradiction. No consistent $\pm 1$
    assignment exists.
\end{enumerate}

\medskip
The equiangular configuration \emph{automatically forces} odd parity
(Section~\ref{sec:characteristic}).
This is not a coincidence --- it is the evaluation of a global
cohomological invariant.

%------------------------------------------------------------------
\section{\texorpdfstring{$[f_3]$}{[f3]} as Symplectic Characteristic Class}
\label{sec:characteristic}
%------------------------------------------------------------------

\begin{corollary}[{${[f_3]}$} as Symplectic Characteristic Class]
\label{cor:characteristic}
The odd parity of every equiangular $K_5$ is the evaluation of
$[f_3] \in H^1(K_5, \cF)$, the symplectic characteristic class of
$W(5, \F_2)$. This is not an ad hoc sign assignment --- it is a global
cohomological invariant of the symplectic building.
\end{corollary}

\begin{proof}[Proof sketch (computational)]
The sheaf $\cF$ assigns to each context $C_i$ the set of valid
$\pm 1$ eigenvalue assignments to its 4 operators, and to each edge
$C_i \cap C_j$ the eigenvalue of the shared operator. The obstruction
class $[f_3] \in H^1(K_5, \cF)$ measures the failure of local sections
to glue into a global section.

By exhaustive computational verification
(Table~\ref{tab:parity-dist}, Script~P8~\cite{PaperXsuppl}), every equiangular $K_5$
in $W(5, \F_2)$ has odd parity $\prod s_i = -1$. Since every proper
context $K_5$ is equiangular (Theorem~\ref{thm:equiangular}), every
Mermin pentagram has $[f_3] \neq 0$.

A purely geometric proof that the equiangular condition forces
$\prod s_i = -1$ --- which would establish $[f_3]$ as a true
symplectic characteristic class without computation --- remains open
(Conjecture~\ref{conj:odd-parity}).
\end{proof}

\begin{remark}[Geometric Interpretation]
The KS obstruction lives in the \emph{symplectic embedding} of $K_5$
into $W(5, \F_2)$, not in individual Lagrangian intersections with $V$.
The equiangular configuration is a rigid geometric structure that forces
the KS parity anomaly via $[f_3]$. Fano collinearity is merely a shadow
of this deeper structure.
\end{remark}

\subsection{Conjecture~2.5: The Wrong Question}

Conjecture~2.5 of Paper~IX~\cite{PaperIX} proposed a collinearity--obstruction
dictionary: the KS obstruction would translate to a collinearity
obstruction in $PG(2, \F_2)$.

This paper shows that Conjecture~2.5 asked the \emph{wrong question}:

\begin{itemize}
  \item Collinearity is a \emph{projection} of the equiangular structure
    onto $V$, not the obstruction itself.
  \item The true invariant is the symplectic embedding of $K_5$ into
    $W(5, \F_2)$: every proper context $K_5$ is equiangular, and this
    rigid geometric structure forces the KS parity anomaly.
  \item The collinearity pattern varies (4-arcs, non-collinear triples,
    zero Fano points), but the equiangular structure is invariant.
\end{itemize}

The correct statement is Theorem~\ref{thm:equiangular}: the equiangular
Lagrangian configuration with its 10-ray cap structure is the true
geometric invariant.

%------------------------------------------------------------------
\section{\texorpdfstring{$G_2(2)$}{G2(2)} Orbit Decomposition}
\label{sec:orbits}
%------------------------------------------------------------------

\begin{proposition}[$G_2(2)$ Orbit Structure, computational]
\label{prop:orbits}
$G_2(2)$ (order $12{,}096$; its simple subgroup $G_2(2)' \cong PSU(3,3)$
has index~2) acts on the 12{,}096 Mermin pentagrams with exactly 4 orbits:
\[
12{,}096 = 6{,}048 + 2{,}016 + 2{,}016 + 2{,}016.
\]
The largest orbit ($\mathcal{O}_1$, size $6{,}048$, stabilizer order~$2$)
contains the standard pentagram. The three equal-sized special orbits
($\mathcal{O}_2, \mathcal{O}_3, \mathcal{O}_4$, size $2{,}016$ each,
stabilizer order~$6$) carry enhanced symmetry.
\end{proposition}

\begin{corollary}[Two Geometric Classes]
\label{cor:two-classes}
The orbit type distributions reveal a splitting into two geometric
classes:

\begin{table}[h]
\centering
\begin{tabular}{l l r r r}
\toprule
Class & Orbits & Size & F2\% & F4\% \\
\midrule
I & $\mathcal{O}_1 \cup \mathcal{O}_2$ & 8{,}064 & 61.1\% & 27.8\% \\
II & $\mathcal{O}_3 \cup \mathcal{O}_4$ & 4{,}032 & 55.6\% & 33.3\% \\
\bottomrule
\end{tabular}
\caption{Two-class structure of Mermin pentagrams.}
\label{tab:two-classes}
\end{table}

Class~II pentagrams have enhanced transverse obstruction ($+5.5$ pp in
Type~F4) and reduced Fano shadow ($-5.5$ pp in Type~F2) relative to
Class~I. The standard pentagram belongs to Class~I.

In all four orbits, the Type~B fraction is exactly $1/9$:
$672/6{,}048 = 224/2{,}016 = 1/9$.
This uniformity suggests that the Type~B pattern (4 Fano points with
one collision) is a $G_2(2)$-invariant.
\end{corollary}

The 3-fold symmetry among $\mathcal{O}_2, \mathcal{O}_3, \mathcal{O}_4$
is partially broken: $\mathcal{O}_2$ shares the type ratios of
$\mathcal{O}_1$, while $\mathcal{O}_3$ and $\mathcal{O}_4$ form a
conjugate pair with distinct ratios.

%------------------------------------------------------------------
\section{Open Problems}
\label{sec:open-problems}
%------------------------------------------------------------------

\subsection{Geometric Proof of Odd Parity}

Theorem~\ref{thm:equiangular} proves that every proper context $K_5$
is equiangular. The computational verification shows that every
equiangular $K_5$ has odd parity (100\%). A geometric proof of this
implication --- without computation --- would establish $[f_3]$ as a
purely cohomological invariant:

\begin{conjecture}
\label{conj:odd-parity}
The equiangular configuration $\dim(L_i \cap L_j) = 1$ for all 10 pairs
implies odd parity $\prod s_i = -1$ by a direct geometric argument.
Equivalently, $[f_3] \in H^1(K_5, \cF)$ is the non-trivial class for
every equiangular $K_5$.
\end{conjecture}

\subsection{Orbit Stabilizer Structure}

Proposition~\ref{prop:orbits} establishes 4 orbits with stabilizer
orders 2 and 6. The stabilizer of order~6 has two candidates ($Z_6$
or $S_3$), and their identification would clarify whether the 3-fold
partial symmetry arises from a normal subgroup action. Separately,
the geometric characterization of Class~II pentagrams
(Corollary~\ref{cor:two-classes}) --- what distinguishes their
enhanced F4 fraction --- is open.

\subsection{Klein Quartic Connection}

The relationship between the equiangular configuration and the Klein
quartic $X(7)$ is open. The 10-ray cap in $PG(5, \F_2)$ may have a
modular interpretation via the $PSL(2, 7)$ action on $X(7)$.

\subsection{\texorpdfstring{$n$-Qubit}{n-Qubit} Generalization}

The equiangular characterization is specific to 3 qubits ($W(5, \F_2)$).
For $n$ qubits, the obstruction ladder continues:
\begin{itemize}
  \item $n=1$: $W(1, \F_2)$, $H^1$, geometric phase
  \item $n=2$: $W(3, \F_2)$, $H^2$, KS contextuality (Peres--Mermin square)
  \item $n=3$: $W(5, \F_2)$, $H^3$, equiangular characterization (this paper)
  \item $n \geq 4$: $W(2n-1, \F_2)$, higher obstructions (open)
\end{itemize}

The modular geometry connection (via hyperbolic embedding) only opens
at $n=3$. Whether the equiangular characterization generalizes to
$n \geq 4$ is open.

\subsection{\texorpdfstring{Gerbe $H^1$}{Gerbe H1} Non-vanishing}

The twisted Penrose transform
$H^1_{\mathcal{G}}(\CP^3, \cO_{\mathcal{G}}(-n-2))$ remains an open
problem from Paper~VIII~\cite{PaperVIII}. The standard Leray analysis
gives $R^1\pi_*\cO_{\mathcal{G}} = 0$, so the $\Z/2$-gerbe obstruction
cannot be detected by classical sheaf cohomology.

%------------------------------------------------------------------
\section*{Appendix: Rigor Self-Assessment}
%------------------------------------------------------------------

\begin{table}[H]
\centering\small
\setlength{\tabcolsep}{4pt}
\begin{tabular}{l c p{6.5cm}}
\toprule
Claim & Rigor & Status \\
\midrule
135 Lagrangians in $W(5, \F_2)$ & $\bullet$ &
  Standard finite geometry; $|PSp(6, \F_2)|/|P_L| = 135$ \\
945 proper contexts & $\bullet$ &
  $135 \times 7$ (complement of each Fano line) \\
12{,}096 Mermin pentagrams & $\bullet$ &
  Exhaustive enumeration (deduped); code in~\cite{PaperXsuppl} \\
Type~A $= 0$ & $\bullet$ &
  Computational verification; geometric proof open \\
Type~B: 33.3\% are 4-arcs & $\bullet$ &
  Computational verification \\
Type~F2: 87.5\% non-collinear & $\bullet$ &
  Computational verification \\
Type~F4: 25\% zero Fano points & $\bullet$ &
  Computational verification \\
Pivot Pentagon Rigidity & $\bullet$ &
  Computational verification; all 4{,}480 have equiangular config \\
Theorem~\ref{thm:equiangular} (Equiangular $\Leftrightarrow$ Mermin) & $\bullet$ &
  \textbf{Geometric proof} (pure algebra, no computation) \\
Theorem~\ref{thm:cap} (10-Ray Cap) & $\bullet$ &
  Computational (Script~P7~\cite{PaperXsuppl}); geometric proof open \\
$[f_3]$ as symplectic characteristic class & $\bullet/\circ$ &
  Cohomological interpretation ($\bullet$); full proof of uniqueness ($\circ$) \\
$G_2(2)$ orbit decomposition & $\bullet$ &
  4 orbits: $6{,}048 + 2{,}016 \times 3$; two geometric classes \\
Klein quartic connection & $\circ$ &
  Open problem \\
$n$-qubit generalization & $\circ$ &
  Open problem \\
\bottomrule
\end{tabular}
\caption{Rigor self-assessment for Paper~X.}
\label{tab:rigor}
\end{table}

%------------------------------------------------------------------
\section*{Acknowledgments}
%------------------------------------------------------------------

The author thanks the mathematical physics reading group for feedback
on the equiangular characterization, and a research collaborator for
identifying the geometric proof of Theorem~\ref{thm:equiangular} and
the interpretation of $[f_3]$ as a symplectic characteristic class.

%------------------------------------------------------------------
\begin{thebibliography}{10}
%------------------------------------------------------------------

\bibitem[Mermin1990]{Mermin1990}
N.~D.~Mermin,
``Simple unified form for the major no-hidden-variables theorems,''
\textit{Am. J. Phys.} \textbf{58}, 731--734 (1990).

\bibitem[PaperVIII]{PaperVIII}
Z.-L.~Chen,
``The $\Phi$ Functor and the $n$-Qubit Obstruction Ladder:
Explicit Transgression and Twisted Penrose Transform,''
\textit{Zenodo, DOI: 10.5281/zenodo.20454120} (2026).

\bibitem[PaperIX]{PaperIX}
Z.-L.~Chen,
``The 3-Qubit Obstruction Ladder: Hyperbolic Embedding,
Mermin Pentagram, and Two Klein Quartic Bridges,''
\textit{Zenodo, DOI: 10.5281/zenodo.20465623} (2026).

\bibitem[PaperXsuppl]{PaperXsuppl}
Z.-L.~Chen,
``The Equiangular Characterization of Mermin Pentagrams:
Computational Verification,''
\textit{Zenodo, DOI: 10.5281/zenodo.20472357} (2026).

\bibitem[LevayPlanatSaniga]{LevayPlanatSaniga}
P.~L\'evay, M.~Planat, and M.~Saniga,
``Grassmannian Connection Between Three- and Four-Qubit Observables,
Mermin's Contextuality and Black Holes,''
\textit{JHEP} \textbf{09} (2013) 037.

\end{thebibliography}

\end{document}
